\section{Analysis}\label{sec:theory}

We state what FOA guarantees, what it does not, and why the remaining trade-off is unavoidable.
Throughout, device $i$ has calibrated centre $\mu^0$, spread $s^0$, threshold $\tau$ and alarm radius
$\rho=\tau s^0$; $t$ counts cycles since the device joined; $L(t)=B_0+Bt$ and $L_s(t)=B_0^s+B^s t$.
FOA keeps four accumulators: the applied common movement $C(t)\in\R^D$ and $C_s(t)\in\R$, and the
applied residual movement $R(t)$ and $R_s(t)$.

\subsection{Containment}
\begin{theorem}[Residual containment]\label{thm:contain}
For any sequence of batches and peer frames, at every cycle
\begin{gather}
\mu(t)=\mu^0+C(t)+R(t),\qquad \log\frac{s(t)}{s^0}=C_s(t)+R_s(t),\\
\norm{R(t)}\le L(t)\,\rho,\qquad |R_s(t)|\le L_s(t),
\end{gather}
and each common increment is $a_t=\kappa_t\,w_t\lambda_t\bar m_t$ with $\kappa_t\in[0,1]$ (the common
budget's shrink factor), where $\bar m_t$ and $w_t$ are functions of the verified peer frames only and
$\lambda_t\in[0,1]$.
\end{theorem}
\begin{proof}
The device changes $\mu$ only by $a_t+r_t$, with $a_t$ added to $C$ and $r_t$ to $R$; by induction from
$\mu(0)=\mu^0$, $C(0)=R(0)=0$ the decomposition holds, and the same for the log-spread. The residual
step sets $R(t)$ to the projection of $R(t-1)+(d_t-a_t)$ onto the ball of radius $L(t)\rho$, so the
bound holds after every step; likewise $R_s$ is clipped to $[-L_s(t),L_s(t)]$. Before the common budget,
$a_t=w_t\lambda_t\bar m_t$; projecting $C+a_t$ onto a ball centred at the origin that contains $C$
returns $C+\kappa_t a_t$ for some $\kappa_t\in[0,1]$. $\bar m_t$ is an average of geometric medians of
peer proposals and $w_t$ a function of their dispersion, so neither reads the device's own traffic.
$\lambda_t$ is clipped to $[0,1]$.
\end{proof}

The theorem says what the victim's own traffic can and cannot do. Its own traffic enters only through
$\lambda_t$, which can only \emph{scale down} a peer-determined vector, and through the residual, which
is bounded by $L(t)\rho$. It cannot add any direction the peers did not vote for. Because $\lambda_t$
uses the median projection of the device's core points, an attacker with less than half the batch can
neither drive it to zero (to stop the device following the fleet) nor to one (to drag the device along a
consensus its own traffic does not show).

\begin{corollary}[Minimum time to compromise]\label{cor:ttc}
Let $d_{50}(t)$ be the median distance, in calibration spreads, of the target attack traffic from the
consensus-backed centre $\mu^0+C(t)$. If at cycle $t$ fewer than half of that traffic is flagged, then
\begin{equation}
d_{50}(t)\;\le\;\tau\, e^{C_s(t)+L_s(t)}+\tau L(t).\label{eq:reach}
\end{equation}
With no common movement ($C=0$, $C_s=0$), compromise is impossible before the first $t$ at which
$d_{50}\le\tau(e^{B_0^s+B^s t}+B_0+Bt)$; for a target at $M$ alarm radii and small budgets this is about
$t^\ast\approx (M-1-B_0-B_0^s)/(B+B^s)$ cycles.
\end{corollary}
\begin{proof}
If more than half the attack points $x$ satisfy $\norm{x-\mu(t)}\le\tau s(t)$, then for those
$\norm{x-\mu^0-C(t)}\le \tau s(t)+\norm{R(t)}$, so the median is at most that. Divide by $s^0$ and apply
Theorem~\ref{thm:contain}. The approximation uses $e^{y}\approx1+y$.
\end{proof}
We checked the bound in every cycle of all 3{,}630 attacked runs of FOA and its budgeted variants and of 120
synthetic runs: it was never violated. The only violations came from the ablation that removes the
residual budget, to which the corollary does not apply (Section~\ref{sec:ablation}), which shows the
bound is not vacuous.

\subsection{Byzantine peers}
\begin{lemma}[Minsker~\cite{minsker}, Lemma 2.1]\label{lem:gm}
Let $0<\alpha<1/2$ and $C_\alpha=(1-\alpha)\sqrt{1/(1-2\alpha)}$. If more than $(1-\alpha)K$ of $K$ points
lie within distance $r$ of $z$, their geometric median lies within $C_\alpha r$ of $z$.
\end{lemma}

\begin{theorem}[Bounded Byzantine influence]\label{thm:byz}
Suppose $f<K/2$ of the $K$ peers hold valid keys and send arbitrary frames, and in cycle $t$ the honest
peers' proposals lie within $r_t$ of their own geometric median $m_t^H$. Then
$\norm{m_t-m_t^H}\le C_{f/K}\,r_t$, so the adversarial part of each common increment is at most
$w_t\lambda_t$ times an average of such terms, and the total adversarial common movement is at most
$\min\big(\sum_u C_{f/K} r_u,\;(B_0^c+B^c t)\rho\big)$.
\end{theorem}
\begin{proof}
Apply Lemma~\ref{lem:gm} with $z=m_t^H$: $K-f>(1-f/K)K$ points lie within $r_t$ of it. The exponential
average is a convex combination, $w_t\lambda_t\le1$, and the common budget caps $\norm{C(t)}$.
\end{proof}
This is a bound on \emph{rate}, not immunity: a minority of compromised peers can still push a
consistent bias of order $C_\alpha r$ per cycle until the common budget stops it, and a majority
($f\ge K/2$) can move the consensus anywhere. Keyless attackers have no effect: frames with a bad tag,
an unknown sender or a stale cycle are dropped before aggregation.

\subsection{The frontier}
\begin{proposition}[Device-only drift and poisoning are indistinguishable]\label{prop:frontier}
Let any update rule see the device's own traffic and its peers' traffic. Consider world~A, in which the
device's benign traffic drifts by a vector $v$ per cycle and its peers do not, and world~B, in which the
device's benign traffic is stationary and an attacker controlling the device's traffic injects
traffic distributed exactly as world~A's drifted traffic. The observations are identically
distributed in both worlds. Hence a rule that keeps false alarms bounded in world~A for every $v$ with
$\norm{v}\le\nu$ moves the detector by up to $\nu$ per cycle in world~B.
\end{proposition}
\begin{proof}
The peers' traffic is identical in A and B and the device's traffic has the same law, so any rule
produces the same distribution of detector trajectories. Tracking A means following the drift.
\end{proof}
The proposition is the multi-stream counterpart of single-stream impossibility results~\cite{warden}:
peers remove \emph{common} drift from the dilemma but cannot remove \emph{device-only} drift. FOA makes
the remaining trade-off a single explicit rate: device-only movement, honest or not, is admitted at
rate $B$ (plus $B^s$ in log-spread). Larger $B$ buys faster tracking of a device's own changes and
proportionally less time to compromise (Corollary~\ref{cor:ttc}); we measure both ends in
Section~\ref{sec:eval}.

\subsection{No freeze state}
\begin{proposition}[Freeze immunity]\label{prop:freeze}
(i) Traffic that the detector alarms on does not enter FOA's update, so injecting such traffic
(decoys) leaves the device's trajectory unchanged. (ii) No input sequence stops FOA from learning: the
residual admits movement of $B\rho$ per cycle forever.
\end{proposition}
\begin{proof}
(i) The local proposal uses only accepted points (score $\le\tau$), and the common part depends only on
peer frames and the median of accepted core points; alarmed points enter neither. (ii) $L(t)$ grows
without bound and nothing in Algorithm~\ref{alg:foa} sets an update to zero based on the device's
alarm rate or any other suspicion statistic.
\end{proof}
Gated defences have exactly the state that Proposition~\ref{prop:freeze} rules out: a condition on a
traffic statistic under which they stop learning. An attacker who can make that condition true by
sending loud traffic can freeze them, and honest drift then turns the freeze into a lasting flood of
false alarms (Section~\ref{sec:eval}).
